\documentclass[15pt]{scrartcl}
\usepackage[sexy]{evan}
\newcommand{\R}{\mathbb{R}}
\newcommand{\N}{\mathbb{N}}
\newcommand{\Q}{\mathbb{Q}}
\newcommand{\Z}{\mathbb{Z}}
\newcommand{\F}{\mathbb{F}}
\usepackage{relsize}
\setcounter{secnumdepth}{0}
\begin{document}
\title{Daily Log on Math and Random Thoughts}
\date{\today}
\author{Jiawen Huang}
\maketitle
This is my daily math log. I update this every other day. This is mainly for my own review and reflection. Do not take any of the contents too seriously. A large number of errors may exist. Email me if you find any mistakes or have any suggestions. 

\newpage
\tableofcontents
\newpage

\section{Acknowledgment}
This is Evan Style LaTeX template. Thanks to Evan Chen for creating such a beautiful style file.
\newpage 
\section{Oct. 28, 2025}
\begin{enumerate}
\item Weyl Algebra $\frac{k<x,y>}{xy-yx-1}$. 
\item Center of Weyl Algebra is $k$ in characteristic 0.
\item In characteristic $p$, the center of Weyl Algebra is $k[x^p,y^p]$.
\item Weyl Algebra is simple.
\item $x$ and $y$ can be seen as operators on $k[t]$ by $x:f\mapsto tf$ and $y:f\mapsto \frac{df}{dt}$. This is amazing. 
I think the structure $xy-yx=1$ can be used to make some algebra olympiad problems. Very rich and interesting structure. 
When you interpret as operators on polynomials, everything becomes very concrete and easy to understand. Perfect for a math olympiad problem.

\end{enumerate}

\section{Oct. 29, 2025}
\begin{enumerate}
\item Splitting field of a polynomial is unique up to isomorphism.
\item Separable. An inseparable polynomial: $x^p-t$ over $\mathbb{F}_p(t)$. $x^p-t=(x-u)^p$ in the algebraic closure. 
\item A polynomial is separable iff its derivative and $f$ are coprime. In particular, over characteristic 0, all irreducible polynomials are separable.
\item A field is perfect if all irreducible polynomials are separable iff every algebraic extension is separable. 
\item A finite extension is separable iff $[F:k]_s=[F:k]$ iff $F=k(\alpha_1,...,\alpha_n)$ where $\alpha_i$ are separable.
\item A finite Normal extension is a splitting field of a polynomial. This is amazing. 
\item Riemann integral is equivalent to Darboux integral.
\end{enumerate}
\section{Oct. 30, 2025}
\begin{enumerate}
\item Riemann integrable iff discontinuities has measure 0.
\item A lot of exercises on separable extensions from Aluffi chapter 7. 
\item Formal power series $k[[x]]$ on math 250 lecture. 
\end{enumerate}
\section{Nov. 1, 2025}
\begin{enumerate}
\item Cyclotomic polynomial is in $\Z[x]$ ,and it's always irreducible. Proof used Dirichlet's Theorem and the fact that finite fields are perfect. Too smart. 
\item Finite Field of order $p^r$ is unique up to isomorphism, it's the splitting field of $x^{p^r}-x$. 
\item $x^4+1$ is reducible in all finite field. But irreducible in $\Z[x]$. Proof use the fact that $x^4+1 \mid x^8-1\mid x^{p^2}-x$. So the root $\beta$ of $x^4+1$ in $F_{p^2}$ has degree of $1$ or $2$. Thus, its minimal polynomial has degree of $1$ or $2$.
\item Let $F=\F_q$, $x^{p^n}-x$ is the product of all the irreducible polynomial of degree $d$ for all $d\mid n$. 
\item The automorphism $Aut_{\F_p}(\F_{p^d})\cong C_q$. The generator is Frobenius function $x\mapsto x^p$. 
\end{enumerate}
\section{Nov. 2, 2025}
\begin{enumerate}
\item There are countably many jump discontinuities. 
\item Set of Critical Value is zero set. 
\item The automorphism $\phi\in Aut_k(k(x))$ must be $\phi: x\mapsto \frac{ax+b}{cx+d}$ where $ac-bd\not=0$. A surprising result but hard to prove. Is there a natural but easy way to see this?
\item Review H104. 
\end{enumerate}
\section{Nov. 3, 2025}
\begin{enumerate}
\item $k\subseteq F$ iff there are finitely many fields $E$ such that $k\subseteq E\subseteq F$. 
This is highly nontrivial to me. My intuition is that if $k(a,b)$ is not a simple extension of $k$, then it may have a lot of possible combination of $a$ and $b$ to form an intermediate field $E$. This turns out to be helpful for one direction. 
\item All finite separable extensions are simple. This is very surprising and makes our life substantially easier. Consequently, it's very rare to see a nonsimple extension. The proof is very hard, and I find it hard to come up with. 
\end{enumerate}
\section{Nov. 4, 2025}
\begin{enumerate}
\item Set up Wolfram. 
\item Weddernburn Little Theorem: all finite division rings are fields. The proof by Ernst Witt is incredibly clever. It uses class formula of group and the cyclotomic polynomial.
\item A Riemann integrable function has integral $0$ iff it's zero almost everywhere (set of roots is a zero set).
\item Antiderivative Theorem: The antiderivative of a Riemann integrable function, if it exists, differs from the indefinite integral by a constant. This can be proved using Mean Value Theorem. 
\end{enumerate}
\section{Nov. 5, 2025}
\begin{enumerate}
\item Zsigmondy's Theorem: For $a>b>0$ coprime integers, $a^n-b^n$ has a prime factor which does not divide $a^k-b^k$ for any $k<n$, with exceptions $(a,b,n)=(2,1,6)$ or $n=2$ and $a+b$ is a power of $2$. The proof used LET lemma and properties of cyclotomic polynomial. Very clever proof.
\item There are amazing interactions with cyclotomic polynomials and order of an element $a$ in the multiplicative group mod $p$ when $p \not\mid n$: $p\mid \Phi_n(a)$ iff the order of $a$ mod $p$ is $n$. Thus, $n\mid p-1$ by Fermat's Little Theorem. With this, we can show that there are infinitely many primes congruent to $1$ mod $n$ by considering $\Phi_n(k)$ for some integer $k$.
\end{enumerate}
\section{Nov. 6, 2025}
\begin{enumerate}
\item Took 250 midterm. All questions are trivial. 
\item Went to an AYCE Sushi at San Francisco.
\item Learn Series in H104, easy. 
\end{enumerate}
\section{Nov. 7, 2025}
\begin{enumerate}
\item Fundamental Theorem of Galois Theory: Let $k\subseteq F$ be a finite Galois extension. The lattice of field extensions between $k$ and $F$ is anti-isomorphic to the lattice of subgroups of $Gal(F/k)$. This is astonishing. The proof is quite straightforward, but I don't know how Galois came up with it. Too deep.  
\item Mrs. Xinyang Qian gives a guest lecture on the examples of a derivative of a function not being Riemann integrable using Cantor fat set and topologist sine curve. Very interesting examples. 
\item Prove that if $f_n$ is a sequence of continuous functions converging uniformly to $f$, then $f$ is continuous. 
\end{enumerate}
\section{Nov. 8, 2025}
\begin{enumerate}
    \item BMT today. I met friends from my high school. Went to AYCE BBQ.
\end{enumerate}
\section{Nov. 9, 2025}
\begin{enumerate}
    \item Finished Analysis homework and algebra homework. Algebra homework was so hard. 
\end{enumerate}
\section{Nov. 10, 2025}
\begin{enumerate}
\item The Second Fundamental Theorem of Galois Theory: $k\subseteq E$ is Galois iff $Gal(F/E)$ is normal in $Gal(F/k)$. In this case, $Gal(E/k)\cong Gal(F/k)/Gal(F/E)$. Proof on Aluffi is very clever: consider the action of $Gal(F/k)$ on the set of embeddings of $E$ into $\overline{k}$, $I$. $I$ as a $Gal(F/k)$-set is isomorphic to the coset space $Gal(F/k)/H$ where $H=Gal(F/E)$. The rest is straightforward.
\item Some random Analysis facts about converging radius and sequence of functions. 
\end{enumerate}
\section{Nov. 11, 2025}
\begin{enumerate}
\item If $k\subseteq F$ is an extension of degree $m$. Assume the characteristic of $k$ doesn't divide $m$ and $k$ contains primitive $m$-th root of unity, 
then $F$ is a Galois Cyclic extension if and only if $F=k(\alpha)$ where $\alpha^m\in k$. 
\item If $\phi_1, \phi_2, \cdots, \phi_n$ are pairwise distinct homomorphisms from a field $F$ to a field $K$, then they are linearly independent over $K$. This can be generalized to integral domains. 
\end{enumerate}
\section{Nov. 12, 2025}
\begin{enumerate}
\item Hilbert's theorem 90: Let $k\subseteq F$ be a cyclic Galois extension with Galois group generated by $\sigma$. 
Then an element $\beta\in F$ has norm $1$ iff there exists an $\alpha\in F$ such that $\beta=\frac{\alpha}{\sigma(\alpha)}$. The proof is very clever. It constructs an element $\theta=\sum_{i=0}^{n-1} \sigma^i(\alpha)x^i$ in $F[x]$ and considers the action of $\sigma$ on it.
\\
This is super powerful. It can be used to solve the Pell's Equations over rationals and other number theoretic problems. 
\item If $k\subseteq F$ and $k\subseteq L$ is Galois, then $k\subseteq FL$ is Galois. I came up with a very clever proof: 
It's clear that $k\subseteq FL$  is finite and separable. 
So we can find its Galois Closure $L$. 
It suffices to show $Aut_{FL}(L)$ is normal in $Aut_k(L)$. 
Notice that $Aut_{FL}(L)=Aut_L(L)\cap Aut_F(L)$. 
Since $k\subseteq L$ and $k\subseteq F$ is Galois, 
$Aut_L(L)$ and $Aut_F(L)$ are normal in $Aut_k(L)$. 
Thus, their intersection $Aut_{FL}(L)$ is normal. 
Thus, $k\subseteq FL$ is Galois.
\end{enumerate}
\section{Nov. 13, 2025}
\begin{enumerate}
\item Watched A Mortal's Journey to Immortality. 
\end{enumerate}
\section{Nov. 14, 2025}
\begin{enumerate}
\item I learned an algebraic proof of the Fundamental Theorem of Algebra using Galois theory. The central idea is that any algebraic extension of $\R$ must be of degree of $2^n$ because polynomials of odd degree must have a real root. Then use Galois theory to convert the problem into a group theory problem.
\item All symmetric functions of $k[x_1,...,x_n]$ are generated by elementary symmetric polynomials. The proof is very clever. Consider the extension $k(e_1,...,e_n)\subseteq k(x_1,...,x_n)$ where $e_i$ are elementary symmetric polynomials. This is a Galois extension with Galois group $S_n$. So the set of symmetric functions is the fixed field of $S_n$, which is $k(e_1,...,e_n)$ by the Fundamental Theorem of Galois Theory.
\item Every group can be realized as a Galois group. The proof is very clever. 
\end{enumerate}
\section{Nov. 15, 2025}
\begin{enumerate}
\item In characteristic $0$, an irreducible polynomial is solvable by radicals iff the splitting field is contained in a Galois radical extension iff the Galois group of the splitting field is a solvable group.
\item Since the Galois group of $k(e_1,...,e_n)\subseteq k(x_1,...,x_n)$, $S_n$, is not solvable for $n\geq 5$, there is no formula using radicals to solve general polynomials of degree at least $5$ because their Galois groups are not solvable. 
\begin{remark}
    This is something I knew since elementary school, but I never really understood why. It's much harder than I thought. Galois is a genius. Today is definitely a remarkable day of my life.
\end{remark}
\item Every finite Abelian group can be realized as a Galois group over $\Q$. The key is cyclotomic extensions and the fact that every finite Abelian group is a homomorphic image of $(\Z/n\Z)^{\times}$ for some $n$. (A weak form of Dirichlet's Theorem is needed here.)
\end{enumerate}

1

\end{document}